\documentclass{article}
\usepackage[utf8]{inputenc}
\usepackage{hyperref}
\usepackage{listings}
\usepackage{xcolor}
\usepackage{enumitem}

\lstset{
    basicstyle=\ttfamily\small,
    breaklines=true,
    frame=single,
    backgroundcolor=\color{gray!5},
}

\title{Kernel errors/warnings summary with full severity levels and colorized output}
\author{Marcos de Carvalho}
\date{2026-04-21}

\begin{document}

\section*{Kernel errors/warnings summary with full severity levels and colorized output}
\textbf{Author:} Marcos de Carvalho \hfill \textbf{Date:} 2026-04-21

\noindent Retrieves all kernel messages from the kernel ring buffer, extracts and counts all message severity levels (emerg, alert, crit, err, warn, notice, info, debug), and presents them with a comprehensive summary frontmatter followed by the human-readable colorized logs.

\section*{Language}
bash

\section*{Category}
diagnostics

\section*{Command}
\begin{lstlisting}
dmesg -T -x --level=emerg,alert,crit,err,warn --color=always 2>/dev/null | awk 'BEGIN { print "\033[1m=== KERNEL MESSAGE SUMMARY ===\033[0m" } { lines[NR] = \$0; lvl = \$2; gsub(/\x1B\[[0-9;]*[mK]/, \"\", lvl); gsub(/[: \t]/, \"\", lvl); if(lvl ~ /^(emerg|alert|crit|err|warn|notice|info|debug)\$/) cnt[lvl]++; if(lvl ~ /^\w+\$/) all_cnt[lvl]++ } END { if(NR>0) { print "\033[1m--- ERRORS & WARNINGS ---\033[0m"; for(l in cnt) printf \"  %-7s : %d\n\", l, cnt[l]; print "\033[1m--- ALL MESSAGES (by severity) ---\033[0m"; for(l in all_cnt) printf \"  %-7s : %d\n\", l, all_cnt[l]; print "\033[1m======================================\033[0m\n"; for(i=1; i<=NR; i++) print lines[i] } else print "No kernel messages found for selected levels." }'
\end{lstlisting}

\section*{Explanation}
The one-liner uses 'dmesg -T -x' to extract kernel messages with human-readable timestamps (-T) and decode the facility/level prefixes (-x). We filter for all severity levels (emerg,alert,crit,err,warn,notice,info,debug) via '--level'. The '--color=always' flag preserves standard color coding. The output is piped into 'awk', which does several tasks: (1) strips ANSI escape sequences from the second column to cleanly count occurrences, (2) maintains two counters: one for 'errors/warnings' (emerg,alert,crit,err,warn) and another for 'all messages' (all severity levels), and (3) stores all log lines in memory. The END block prints a formatted summary frontmatter with both the 'errors/warnings' and 'all messages' counts, followed by the original colorized logs.

\section*{Tags}
dmesg, kernel, errors, warnings, logs, troubleshooting, monitoring, summary, awk, diagnostics, severity, log-analysis, linux, system, administration

\section*{Dependencies}
util-linux, gawk


\section*{Arguments}
None

\section*{Examples}
\begin{enumerate}
  \item \texttt{dmesg -T -x --level=emerg,alert,crit,err,warn --color=always 2>/dev/null | awk 'BEGIN \{ print "\textbackslash\{\}033[1m=== KERNEL MESSAGE SUMMARY ===\textbackslash\{\}033[0m" \} \{ lines[NR] = \textbackslash\{\}\$0; lvl = \textbackslash\{\}\$2; gsub(/\textbackslash\{\}x1B\textbackslash\{\}[[0-9;]*[mK]/, \textbackslash\{\}"\textbackslash\{\}", lvl); gsub(/[: \textbackslash\{\}t]/, \textbackslash\{\}"\textbackslash\{\}", lvl); if(lvl \textasciitilde{} /\textasciicircum{}(emerg|alert|crit|err|warn|notice|info|debug)\textbackslash\{\}\$/) cnt[lvl]++; if(lvl \textasciitilde{} /\textasciicircum{}\textbackslash\{\}w+\textbackslash\{\}\$/) all\_cnt[lvl]++ \} END \{ if(NR>0) \{ print "\textbackslash\{\}033[1m--- ERRORS \& WARNINGS ---\textbackslash\{\}033[0m"; for(l in cnt) printf \textbackslash\{\}"  \%-7s : \%d\textbackslash\{\}n\textbackslash\{\}", l, cnt[l]; print "\textbackslash\{\}033[1m--- ALL MESSAGES (by severity) ---\textbackslash\{\}033[0m"; for(l in all\_cnt) printf \textbackslash\{\}"  \%-7s : \%d\textbackslash\{\}n\textbackslash\{\}", l, all\_cnt[l]; print "\textbackslash\{\}033[1m======================================\textbackslash\{\}033[0m\textbackslash\{\}n"; for(i=1; i<=NR; i++) print lines[i] \} else print "No kernel messages found for selected levels." \}'} - Show full severity summary and colorized kernel messages
  \item \texttt{dmesg -T -x --level=emerg,alert,crit,err,warn --color=always 2>/dev/null | awk 'BEGIN \{ print "\textbackslash\{\}033[1m=== KERNEL MESSAGE SUMMARY ===\textbackslash\{\}033[0m" \} \{ lines[NR] = \textbackslash\{\}\$0; lvl = \textbackslash\{\}\$2; gsub(/\textbackslash\{\}x1B\textbackslash\{\}[[0-9;]*[mK]/, \textbackslash\{\}"\textbackslash\{\}", lvl); gsub(/[: \textbackslash\{\}t]/, \textbackslash\{\}"\textbackslash\{\}", lvl); if(lvl \textasciitilde{} /\textasciicircum{}(emerg|alert|crit|err|warn|notice|info|debug)\textbackslash\{\}\$/) cnt[lvl]++; if(lvl \textasciitilde{} /\textasciicircum{}\textbackslash\{\}w+\textbackslash\{\}\$/) all\_cnt[lvl]++ \} END \{ if(NR>0) \{ print "\textbackslash\{\}033[1m--- ERRORS \& WARNINGS ---\textbackslash\{\}033[0m"; for(l in cnt) printf \textbackslash\{\}"  \%-7s : \%d\textbackslash\{\}n\textbackslash\{\}", l, cnt[l]; print "\textbackslash\{\}033[1m--- ALL MESSAGES (by severity) ---\textbackslash\{\}033[0m"; for(l in all\_cnt) printf \textbackslash\{\}"  \%-7s : \%d\textbackslash\{\}n\textbackslash\{\}", l, all\_cnt[l]; print "\textbackslash\{\}033[1m======================================\textbackslash\{\}033[0m\textbackslash\{\}n"; for(i=1; i<=NR; i++) print lines[i] \} else print "No kernel messages found for selected levels." \}' | less -R} - Show the summary and logs with a pager, preserving colors
  \item \texttt{sudo dmesg -T -x --level=emerg,alert,crit,err,warn --color=always 2>/dev/null | awk 'BEGIN \{ print "\textbackslash\{\}033[1m=== KERNEL MESSAGE SUMMARY ===\textbackslash\{\}033[0m" \} \{ lines[NR] = \textbackslash\{\}\$0; lvl = \textbackslash\{\}\$2; gsub(/\textbackslash\{\}x1B\textbackslash\{\}[[0-9;]*[mK]/, \textbackslash\{\}"\textbackslash\{\}", lvl); gsub(/[: \textbackslash\{\}t]/, \textbackslash\{\}"\textbackslash\{\}", lvl); if(lvl \textasciitilde{} /\textasciicircum{}(emerg|alert|crit|err|warn|notice|info|debug)\textbackslash\{\}\$/) cnt[lvl]++; if(lvl \textasciitilde{} /\textasciicircum{}\textbackslash\{\}w+\textbackslash\{\}\$/) all\_cnt[lvl]++ \} END \{ if(NR>0) \{ print "\textbackslash\{\}033[1m--- ERRORS \& WARNINGS ---\textbackslash\{\}033[0m"; for(l in cnt) printf \textbackslash\{\}"  \%-7s : \%d\textbackslash\{\}n\textbackslash\{\}", l, cnt[l]; print "\textbackslash\{\}033[1m--- ALL MESSAGES (by severity) ---\textbackslash\{\}033[0m"; for(l in all\_cnt) printf \textbackslash\{\}"  \%-7s : \%d\textbackslash\{\}n\textbackslash\{\}", l, all\_cnt[l]; print "\textbackslash\{\}033[1m======================================\textbackslash\{\}033[0m\textbackslash\{\}n"; for(i=1; i<=NR; i++) print lines[i] \} else print "No kernel messages found for selected levels." \}'} - Run with sudo to access kernel ring buffer on restricted systems
\end{enumerate}

\section*{Output}
\begin{lstlisting}
\u001b[1m=== KERNEL MESSAGE SUMMARY ===\u001b[0m
\u001b[1m--- ERRORS & WARNINGS ---\u001b[0m
  err     : 2
  warn    : 3
  crit    : 1
\u001b[1m--- ALL MESSAGES (by severity) ---\u001b[0m
  err     : 2
  warn    : 3
  crit    : 1
  notice  : 5
  info    : 12
  debug   : 0
\u001b[1m======================================\u001b[0m

kern  :err   : [Tue Apr 21 10:00:00 2026] sd 0:0:0:0: [sda] Write cache: enabled
kern  :warn  : [Tue Apr 21 10:01:23 2026] WARNING: CPU: 0 PID: 1234
kern  :crit  : [Tue Apr 21 10:02:45 2026] Critical hardware error
kern  :warn  : [Tue Apr 21 10:05:00 2026] ext4-fs: mounted with warnings
kern  :err   : [Tue Apr 21 10:06:12 2026] usb 1-1: device descriptor read/64, error -71
kern  :warn  : [Tue Apr 21 10:08:34 2026] IPv6: ADDRCONF(NETDEV_UP): eth0: link is not ready
kern  :notice: [Tue Apr 21 10:10:00 2026] Kernel log system initialized
kern  :info  : [Tue Apr 21 10:15:00 2026] Network interface eth0 up and running
\end{lstlisting}

\section*{Notes}
\begin{itemize}
  \item Requires read access to the kernel ring buffer, which often requires sudo or being in the 'adm' group on some distributions.
  \item The '-x' (--decode) flag is critical here as it explicitly outputs the severity level on every line, making awk parsing reliable.
  \item ANSI stripping regex in awk ensures color codes don't interfere with severity counting.
  \item This command counts all message severity levels: emerg, alert, crit, err, warn, notice, info, debug.
  \item The 'errors \& warnings' summary includes: emerg, alert, crit, err, warn.
  \item The 'all messages' summary includes all severity levels, providing complete log statistics.
  \item The entire dmesg output is stored in awk's memory (the 'lines' array) before printing, which is perfectly safe for typical volumes but could consume slightly more memory on massively flooded systems.
\end{itemize}

\section*{Warnings}
\begin{itemize}
  \item Running without sufficient privileges (like sudo) will result in a blank output or a permission denied error on most modern secure Linux distributions.
\end{itemize}

\section*{See Also}
\begin{itemize}
  \item find-large-files-recursive
  \item disk-space-sort-largest-directories
\end{itemize}

\section*{Status}
reviewed

\section*{Safety}
safe

\section*{Shell}
bash

\section*{Platforms}
linux

\section*{Created At}
2026-04-21

\section*{Updated At}
2026-04-21

\end{document}
